\begin{table}[htbp]\centering\small
\setlength{\tabcolsep}{4pt}
\begin{tabular}{r r r r r r}
\hline
$n$ & $N_n$ & $d_n$ & $|F_n(v_n)|$ & $\Delta V_n/H_H^2$ & $Q_n^{\rm eff}$ \\
\hline
1 & 0.69315 & 4 & 1.41992 & 0.427528 & 24.350 \\
2 & 1.09861 & 9 & 0.94168 & 0.096593 & 9.405 \\
3 & 1.38629 & 16 & 0.70426 & 0.060098 & 8.247 \\
4 & 1.60944 & 25 & 0.56259 & 0.044692 & 7.907 \\
5 & 1.79176 & 36 & 0.46843 & 0.035895 & 7.772 \\
6 & 1.94591 & 49 & 0.40131 & 0.030117 & 7.713 \\
7 & 2.07944 & 64 & 0.35103 & 0.025998 & 7.686 \\
8 & 2.19722 & 81 & 0.31195 & 0.022899 & 7.675 \\
\hline
\end{tabular}
\caption{Discrete $S^3$ shell crossings. Harmonic $n$ enters the long-wavelength sector at $N_n=\ln(n+1)$ exactly, injecting variance $\Delta V_n$. The effective noise amplitude $Q_n^{\rm eff}$ falls steeply and crosses the continuum value $Q_\nu(1)=7.8686$ between the fourth and fifth shells, undershooting it by $2.5\%$ at $n=8$; the first shell exceeds it by a factor of three.}
\label{tab:shells}
\end{table}
